%%%%%%%%%%%%%%%%%%%%%%%%%%%%%%%%%
%
%       EXERCÍCIO
%
%%%%%%%%%%%%%%%%%%%%%%%%%%%%%%%%%

\ifdefstring{\atividade}{prova}{%
    \ifdefstring{\modo}{objetivo}{%
        \renewcommand{\valorquestao}{\ValorQObj\ ponto}
    }{%
        \renewcommand{\valorquestao}{\ValorQDisc\ pontos}
    }%
}%

\begin{exercicioBanco}[\valorquestao]
Em uma cidade, dois aplicativos de transporte, $A$ e $B$, utilizam modelos distintos para calcular o preço de uma corrida em função da distância percorrida.

O aplicativo $A$ cobra uma taxa fixa de R\$ 6,00 mais R\$ 2,00 por quilômetro rodado. Já o aplicativo $B$ cobra uma taxa fixa de R\$ 10,00 mais R\$ 1,50 por quilômetro rodado.

Seja $x$ a distância, em quilômetros, de uma corrida. Determine para quais valores de $x$, em quilômetros, o aplicativo $A$ é mais vantajoso que o aplicativo $B$.

% Define as alternativas
\newcommand{\alternativas}{%
    \begin{center}
        \begin{tabularx}{\textwidth}{XXX}
            (a) \([8,\infty[\). &
            (b) \(]8,\infty[\). &
            (c) \([0,8]\). \\[5pt]
            (d) \([0,8[\). &
            (e) \([0, \infty[\).
        \end{tabularx}
    \end{center}
}

% Resposta correta
\newcommand{\resposta}{D}

% Lógica condicional para exibição
\ifdefstring{\atividade}{lista}{%
    \alternativas
    \vspace{0.5em}
    
    \noindent\textbf{Resposta:} letra \textbf{\resposta}.
}{%
    \ifdefstring{\modo}{objetivo}{%
        \alternativas
    }{%
        % modo = discursiva → não mostra alternativas
    }
}
\end{exercicioBanco}

